\section{Conclusions}

\begin{frame}{Conclusions}
	In this course:
	\begin{itemize}
		\item We reviewed the \textbf{history} of neural networks.
		\pause
		\item We discussed the first mathematical neuron, the \textbf{Perceptron}, and its structural limitations.
		\pause
		\item We explored how to overcome these limitations by combining \textbf{multiple neurons} using \textbf{non-linear activation functions}.
		\pause
		\item We then saw that a new \textbf{learning algorithm} was required.
	\end{itemize}
\end{frame}

\begin{frame}{Conclusions}
	\begin{itemize}
		\item We explored how \textbf{gradient backpropagation} routes gradients through the various operators of a dense (or fully connected) network.
		\pause
		\item Using the computed gradients, we covered different algorithms to \textbf{update parameters} throughout training using batch methods.
		\pause
		\item Finally, we discussed practical solutions to properly \textbf{initialize} neural network parameters.
	\end{itemize}
\end{frame}